\documentclass[11pt,letterpaper]{article}
\usepackage[margin=0.78in]{geometry}
\usepackage{graphicx}
\usepackage{booktabs}
\usepackage{tabularx}
\usepackage{longtable}
\usepackage{array}
\usepackage{enumitem}
\usepackage{xcolor}
\usepackage{fancyhdr}
\usepackage{titlesec}
\usepackage{hyperref}
\usepackage{listings}
\usepackage{newtxtext,newtxmath}
\usepackage{microtype}
\definecolor{navy}{HTML}{173B64}
\definecolor{teal}{HTML}{087E8B}
\definecolor{softblue}{HTML}{F1F6FA}
\hypersetup{colorlinks=true,linkcolor=navy,urlcolor=teal}
\setlength{\parindent}{0pt}
\setlength{\parskip}{5pt}
\setlist[itemize]{leftmargin=5mm,itemsep=2pt,topsep=3pt}
\setlist[enumerate]{leftmargin=6mm,itemsep=2pt,topsep=3pt}
\renewcommand{\arraystretch}{1.18}
\pagestyle{fancy}
\fancyhf{}
\fancyhead[L]{AI Legal Document Intelligence Platform}
\fancyhead[R]{Detailed Viva Study Material}
\fancyfoot[C]{\thepage}
\setlength{\headheight}{14pt}
\titleformat{\section}{\Large\bfseries\color{navy}}{}{0pt}{}
\titleformat{\subsection}{\large\bfseries\color{teal}}{}{0pt}{}
\titleformat{\subsubsection}{\normalsize\bfseries}{}{0pt}{}
\titlespacing*{\section}{0pt}{10pt}{5pt}
\titlespacing*{\subsection}{0pt}{8pt}{4pt}
\lstset{basicstyle=\ttfamily\small,breaklines=true,frame=single,backgroundcolor=\color{softblue},columns=fullflexible,showstringspaces=false}

\begin{document}
\begin{titlepage}
\centering
\vspace*{0.5cm}
\includegraphics[height=4.7cm]{assets/source-image-000.png}\par
\vspace{0.8cm}
{\Huge\bfseries\color{navy} AI Legal Document Intelligence Platform\par}
\vspace{0.45cm}
{\LARGE\bfseries Detailed Mid-Term Viva Study Material\par}
\vspace{0.35cm}
{\large Theory, implementation, research basis and evaluation\par}
\vspace{1cm}
\begin{tabularx}{0.92\textwidth}{>{\bfseries}p{0.16\textwidth}X}
\toprule
Member & Material ownership\\
\midrule
Puru & Starts the presentation: project introduction, problem statement, complete system architecture, end-to-end data flow, technology choices, APIs and the meaning of ``model'' in the present system.\\
Sujal & Research foundation: RAG, legal NLP, graph retrieval, evidence grounding, literature survey, research gaps, proposed contribution and research evaluation logic.\\
Aniket & Engineering detail: repository and folder structure, current implementation, data structures, tests, preliminary findings, limitations, security and remaining eight-month roadmap.\\
\bottomrule
\end{tabularx}
\vfill
{\large Project Supervisor: Mr. Piyush Kushwaha\par}
{\large September 2026\par}
\end{titlepage}

\tableofcontents
\newpage

\section{Shared Foundation: Every Member Must Understand}
\subsection{Project definition}
The project is an evidence-first legal document intelligence platform for contract review. A user uploads a contract PDF and asks a question. The system extracts clauses, preserves page numbers, detects explicit references, retrieves relevant clauses, follows their dependency links, checks whether required evidence is missing, and then either produces a cited answer or refuses to infer absent content.

It is a decision-support prototype, not legal advice and not a replacement for a qualified lawyer. Its technical goal is to make contract question answering more complete, traceable and cautious.

\subsection{Why contracts are difficult for ordinary search}
A contract is not simply a sequence of independent paragraphs. Its meaning is distributed across a hierarchy:
\begin{itemize}
  \item a definition may control the meaning of a word used many pages later;
  \item an exception may narrow a general obligation;
  \item a termination clause may alter the duration of another clause;
  \item a schedule or exhibit may contain the actual values, assets or services;
  \item an amendment may replace an earlier term.
\end{itemize}

Keyword search can locate repeated words, but it does not prove that all legally connected evidence was retrieved. Fixed-size token chunking can cut a clause in the middle and separate a heading, condition or exception from its body. A generator may then write a fluent answer from incomplete evidence.

\subsection{The central design principle}
The project separates deterministic evidence operations from language generation. Parsing, clause identity, cross-reference resolution, graph traversal, missing-reference detection and citation validation are controlled by code. Gemini is optional and is used only after evidence has been selected and verified.

\subsection{What the project is called technically}
The best description is:
\begin{quote}
\textbf{A structure-aware, dependency-verified retrieval system with optional evidence-constrained language generation.}
\end{quote}
It is related to retrieval-augmented generation, but the current baseline is not a full vector RAG system. The current seed retriever is lexical, graph expansion supplies connected evidence, and a deterministic verifier can block generation.

\subsection{Current implementation versus future work}
\begin{tabularx}{\textwidth}{X X}
\toprule
\textbf{Implemented in the current prototype} & \textbf{Planned or incomplete}\\
\midrule
Real PDF upload in Live Services Mode & Dense embeddings and BM25 hybrid retrieval\\
PyMuPDF text and page extraction & Benchmark-scale quantitative evaluation\\
Rule-based legal clause segmentation & OCR and advanced layout analysis\\
Cross-reference normalization and graph edges & Multi-document amendment precedence\\
Lexical retrieval and two-hop traversal & Full PostgreSQL/pgvector production operation\\
Missing-evidence abstention & Authentication and user workspaces\\
Optional Gemini generation and citation validation & Retention, deletion, encryption and audit controls\\
Evidence interface and offline demo mode & Production deployment and monitoring\\
\bottomrule
\end{tabularx}

\newpage
\section{Puru's Study Material: Introduction and Architecture}
\subsection{Problem statement}
The project addresses the risk of answering contract questions from incomplete or structurally fragmented evidence. The core problem is not only finding a similar paragraph. The system must determine what legal unit the paragraph belongs to, which other clauses it references, whether those targets exist, and how the evidence should be presented to a reviewer.

The practical questions are:
\begin{enumerate}
  \item Can the PDF be converted into stable, page-aware legal clauses?
  \item Can explicit legal dependencies be represented and followed?
  \item Can the system detect when a referenced attachment or clause is absent?
  \item Can a generated answer be restricted to the verified evidence?
  \item Can a human inspect the complete evidence path?
\end{enumerate}

\subsection{System objectives}
\begin{itemize}
  \item Preserve legal structure instead of cutting text at arbitrary token counts.
  \item Retain provenance: clause ID, heading, page number and source contract.
  \item Represent references as directed relations between clauses.
  \item Expand retrieval over those relations for up to two hops in the prototype.
  \item Identify unresolved references relevant to the user's question.
  \item Abstain before generation when decisive evidence is absent.
  \item Show direct evidence, linked evidence, missing references and citations in the UI.
\end{itemize}

\subsection{Architecture layers}
\begin{longtable}{>{\bfseries}p{0.23\textwidth}p{0.70\textwidth}}
\toprule
Layer & Responsibility\\
\midrule
React + TypeScript frontend & Accepts file uploads and questions, selects live or offline mode, displays clauses, page numbers, dependency hops, provider information and abstention reasons.\\
Node.js + Express API & Main orchestration layer. Validates uploads, calls Python, stores results, executes retrieval and graph traversal, performs evidence verification, calls Gemini and returns one response to the frontend.\\
Python + FastAPI parser & Dedicated document-processing service. It receives PDF bytes, uses PyMuPDF, identifies legal headings, produces clause objects and extracts cross-references.\\
Storage layer & PostgreSQL schema contains contracts, clauses and clause dependencies. The current viva setup can use an in-memory fallback when PostgreSQL is unavailable.\\
Retrieval and graph service & Scores clauses lexically, chooses seed clauses, follows outgoing dependencies, records hop levels and detects unresolved references.\\
Gemini service & Converts verified evidence into concise JSON with an answer and clause citations. It is bypassed during an abstention and has an evidence-only fallback when no key is configured.\\
\bottomrule
\end{longtable}

\subsection{End-to-end ingestion flow}
\begin{enumerate}
  \item The browser places the selected PDF in a multipart form under the field name \texttt{file}.
  \item React sends it to \texttt{POST /api/contracts/upload} on the Node service.
  \item Multer holds the file in memory, checks that it is a PDF and enforces a 50 MB limit.
  \item Node sends the same binary data to the FastAPI \texttt{/parse-pdf} endpoint.
  \item PyMuPDF opens the PDF and reads text blocks page by page.
  \item The parser detects Article, Section, Clause, Recitals, Schedule, Exhibit, Appendix and Annex headings using regular expressions.
  \item Text following a heading is accumulated until the next heading. The result becomes one clause object.
  \item Cross-reference patterns find expressions such as Section 1.1, Article II, Schedule A and Exhibit B.
  \item Node stores the contract and clauses. It converts reference strings to canonical IDs and creates dependency rows.
  \item The API returns the contract ID, clause count, linked dependency count, processing time and extracted clauses to the interface.
\end{enumerate}

\subsection{End-to-end query flow}
\begin{enumerate}
  \item React sends the question, contract ID, \texttt{top\_k} and hop limit to \texttt{POST /api/query}.
  \item The backend tokenizes the question and calculates a lexical score for every clause.
  \item The highest scoring clauses become Hop 0 seeds.
  \item The graph traversal follows outgoing edges from each seed. Referenced clauses become Hop 1, and their referenced clauses can become Hop 2.
  \item A visited set prevents the same clause from being returned repeatedly and protects against cycles.
  \item The verifier inspects unresolved dependency edges connected to retrieved clauses and checks whether they are relevant to the question.
  \item If a relevant reference is missing, \texttt{is\_abstain} becomes true and Node returns a deterministic insufficient-evidence response.
  \item Otherwise, direct and expanded clauses form the evidence context for Gemini or the evidence-only fallback.
  \item The frontend renders the four-stage chain: seed retrieval, graph traversal, reference verification and final evidence or refusal.
\end{enumerate}

\newpage
\subsection{What ``model'' means in this project}
The word model can cause confusion. The current prototype does not train a new neural network. It combines several models or algorithms:

\begin{tabularx}{\linewidth}{>{\bfseries\raggedright\arraybackslash}p{0.22\linewidth}>{\raggedright\arraybackslash}X}
\toprule
Type & Meaning in this project\\
\midrule
Document model & A contract is represented as clauses with IDs, headings, pages, content and reference lists.\\
Graph model & Each clause is a node. A legal reference creates a directed edge from the source clause to the target clause.\\
Retrieval model & A deterministic lexical scoring rule ranks clauses. This is the current baseline, not a trained embedding model.\\
Verification model & Rule-based logic checks whether references from relevant clauses resolve to real nodes.\\
Generative model & Gemini is an external pretrained large language model. It is not trained by the project; it receives only verified evidence.\\
Future retrieval model & Dense embeddings and hybrid ranking will be evaluated later against the lexical baseline.\\
\bottomrule
\end{tabularx}

\subsection{Lexical scoring in the current code}
Question terms are matched against the clause ID, section title and content. The current code gives stronger weight to structural fields:
\begin{itemize}
  \item a title match adds 5 points;
  \item a clause-ID match adds 4 points;
  \item content occurrences add up to 3 points per term;
  \item if the question explicitly names a cross-reference stored in a clause, 20 points are added.
\end{itemize}
This is simple and explainable, which makes it a useful mid-term baseline. Its weakness is vocabulary mismatch: a semantically relevant clause may use different words. Hybrid semantic retrieval is therefore part of future work.

\subsection{Clause dependency graph}
Suppose clause SEC-1.2 states that obligations under Section 1.1 do not apply to material in Schedule A. The graph contains nodes SEC-1.2, SEC-1.1 and SCHED-A, with directed edges SEC-1.2 $\rightarrow$ SEC-1.1 and SEC-1.2 $\rightarrow$ SCHED-A.

If the parser sees ``Schedule B'' but no SCHED-B node exists, the dependency row is still stored with a null target. This unresolved edge is the evidence that the attachment is missing.

\subsection{Gemini control and citation validation}
The backend constructs an evidence block in this form:
\begin{lstlisting}
[SEC-1.2] Page 1 - Exclusions from Protection
<clause text>

[SEC-1.1] Page 1 - Confidential Information
<linked clause text>
\end{lstlisting}
Gemini is instructed to use only that evidence, treat the contract text as untrusted data, avoid legal advice, return JSON, and cite only supplied clause IDs. The backend creates an allow-list of evidence IDs and rejects any returned citation outside that set. A low temperature is used to reduce variation. This does not prove that every sentence is correct, but it makes the response more controlled and auditable.

\subsection{Technology-choice reasoning}
\begin{itemize}
  \item \textbf{React and TypeScript:} component-based UI plus type checking for structured contract and query responses.
  \item \textbf{Node.js and Express:} suitable for web APIs, JSON, multipart coordination and integration with React.
  \item \textbf{Python and FastAPI:} strong PDF/NLP ecosystem and concise typed APIs.
  \item \textbf{PyMuPDF:} fast extraction of page-aware text blocks.
  \item \textbf{PostgreSQL:} relational storage for contracts, clauses and auditable dependency edges.
  \item \textbf{pgvector:} planned storage for embeddings when semantic retrieval is added.
  \item \textbf{Redis/BullMQ:} optional background processing for expensive embedding and enrichment jobs.
\end{itemize}

\newpage
\section{Sujal's Study Material: Research Foundation and Gaps}
\subsection{Retrieval-Augmented Generation}
RAG combines a retriever with a generator. The retriever selects external passages related to a question, and the generator conditions its response on those passages. This is useful for private contracts because the document is not expected to be part of the language model's training data.

Retrieval does not automatically guarantee evidence completeness. A retriever can miss a referenced clause, and a generator can ignore or misstate retrieved evidence. The project therefore adds structure preservation, dependency expansion, deterministic verification and citation validation around generation.

\subsection{Structural chunking versus fixed-token chunking}
Fixed-token chunking divides text after a selected number of tokens, sometimes with overlap. It is simple and works across domains, but a boundary can occur inside a legal clause. Structural chunking instead uses the document's own headings and hierarchy.
\par\smallskip
\begin{longtable}{>{\bfseries\raggedright\arraybackslash}p{0.17\textwidth}>{\raggedright\arraybackslash}p{0.29\textwidth}>{\raggedright\arraybackslash}p{0.34\textwidth}}
\toprule
Property & Fixed-token chunks & Structural legal clauses\\
\midrule
Boundary & Token count & Article, section, clause or schedule heading\\
Meaning & May be partial & Intended to remain legally coherent\\
Citation & Approximate passage & Clause ID and page number\\
Reference handling & Usually plain text & Extracted into explicit dependency fields\\
Weakness & Fragmentation & Depends on heading and layout recognition\\
\bottomrule
\end{longtable}

\subsection{Graph-based retrieval theory}
A graph is $G=(V,E)$, where $V$ is the set of clause nodes and $E$ is the set of directed reference edges. Seed retrieval produces a subset $S_0 \subseteq V$. Multi-hop expansion produces:
\[
S_{h+1}=\{v \in V \mid \exists u \in S_h, (u,v) \in E\}.
\]
The final evidence set is the union of the seed and expanded nodes up to the configured hop limit. Two hops capture nearby dependencies while limiting context growth. A visited set prevents duplicates and cycles.

\subsection{Missing-evidence verification}
An edge may contain a reference code even when it has no resolved target node. For a retrieved source clause $u$, an unresolved edge $(u,r)$ means that the contract mentions reference $r$, but the parser did not find its target. If the question is about $r$, the system returns an abstention. This is stronger than asking the LLM to guess whether a schedule exists because it is based on the parsed document inventory.

\subsection{Evidence-constrained generation}
Generation is the last stage, not the retrieval mechanism. The evidence prompt contains only selected clauses and explicit rules. Output is requested as structured JSON. Citations must belong to the backend allow-list. A generation failure does not destroy retrieval results; the system can still return the evidence for human review.

\subsection{Primary papers and how the gaps are derived}
The gaps are a synthesis. The papers do not state our exact proposed gap word for word. We compare what each work solves with what contract evidence review still requires.
\begin{longtable}{>{\raggedright\arraybackslash}p{0.18\textwidth}>{\raggedright\arraybackslash}p{0.27\textwidth}>{\raggedright\arraybackslash}p{0.43\textwidth}}
\toprule
\textbf{Paper} & \textbf{Main contribution} & \textbf{Gap or project connection}\\
\midrule
\endfirsthead
\toprule
\textbf{Paper} & \textbf{Main contribution} & \textbf{Gap or project connection}\\
\midrule
\endhead
Lewis et al., 2020, RAG & Combines generation with retrieved non-parametric memory for knowledge-intensive NLP. & Establishes retrieval before generation. It is general-purpose and does not itself guarantee legal-clause preservation, cross-reference completeness or missing-attachment abstention.\\
\addlinespace
Hendrycks et al., 2021, CUAD & More than 13,000 expert annotations for important contract-review clauses. Transformer performance remains challenging and depends on model and data design. & Shows that contract understanding requires specialist data and expert-defined categories. CUAD locates review-relevant text; our question concerns dependency-aware retrieval and verified evidence for QA.\\
\addlinespace
Guha et al., 2023, LegalBench & A benchmark of 162 tasks covering six forms of legal reasoning, built with legal professionals. & Demonstrates that legal reasoning is not one generic task. Our narrower focus is contract evidence retrieval, citation and refusal behavior.\\
\addlinespace
Edge et al., 2024, GraphRAG & Builds an entity graph and community summaries to answer global questions over large private corpora. & Supports relationships beyond isolated passages. Its graph targets entity summarization; ours uses explicit legal clause references and unresolved attachments for deterministic verification.\\
\bottomrule
\end{longtable}

\subsection{Direct paper links}
\begin{itemize}
  \item Lewis et al., \textit{Retrieval-Augmented Generation for Knowledge-Intensive NLP Tasks}: \url{https://arxiv.org/abs/2005.11401}
  \item Hendrycks et al., \textit{CUAD: An Expert-Annotated NLP Dataset for Legal Contract Review}: \url{https://arxiv.org/abs/2103.06268}
  \item Guha et al., \textit{LegalBench: A Collaboratively Built Benchmark for Measuring Legal Reasoning in Large Language Models}: \url{https://arxiv.org/abs/2308.11462}
  \item Edge et al., \textit{From Local to Global: A Graph RAG Approach to Query-Focused Summarization}: \url{https://arxiv.org/abs/2404.16130}
\end{itemize}

\newpage
\subsection{The five research gaps}
\subsubsection{1. Semantic fragmentation}
Generic chunking can divide a clause or separate its heading from its conditions. The hypothesis is that clause-aware segmentation improves evidence quality and citation clarity. The limitation is that rule-based headings may fail on unusual formatting or scanned documents.

\subsubsection{2. Cross-clause dependency loss}
Top-$k$ similarity retrieval returns individually similar passages but may omit a definition or schedule referenced by the best passage. The project explicitly stores and traverses those references.

\subsubsection{3. Missing-evidence hallucination}
If a source clause refers to an absent schedule, the correct response is not to infer its content. The project converts this absence into a testable condition and a refusal result.

\subsubsection{4. Opaque reasoning and citation drift}
A generated answer may not reveal its evidence, and citations may refer to unsupported passages. The project shows evidence stages and validates generated clause IDs.

\subsubsection{5. Multi-document conflict and supersession}
An amendment can override an older agreement, while standard retrieval may return both. The current prototype does not solve this. Future work will model document identity, dates, amendment links and precedence relations.

\subsection{Research question and hypotheses}
\textbf{Research question:} Can structure-aware clause retrieval with dependency verification improve citation correctness and missing-evidence refusal accuracy compared with fixed-size RAG?
\begin{itemize}
  \item H1: structural chunks achieve higher evidence precision than fixed 500-token chunks;
  \item H2: graph expansion improves recall for questions requiring referenced clauses;
  \item H3: deterministic unresolved-reference checks improve refusal accuracy;
  \item H4: citation allow-list validation reduces unsupported citations;
  \item H5: these controls introduce acceptable latency for interactive review.
\end{itemize}

\subsection{Evaluation design}
Build a question set with single-clause, cross-clause and missing-attachment questions. Run the same documents and questions through a fixed-token baseline and the structural graph pipeline. Human-labelled gold evidence should specify which clauses are necessary and whether the answer should be refused.

Metrics include clause-boundary precision, recall and F1; retrieval recall@$k$; mean reciprocal rank; evidence precision; dependency-path recall; citation validity and completeness; answer correctness; abstention precision and recall; false-refusal rate; and ingestion/query latency.

\newpage
\section{Aniket's Study Material: Repository, Testing and Roadmap}
\subsection{Repository structure}
\begin{lstlisting}
major project/
|-- frontend/
|   |-- src/App.tsx
|   |-- src/services/api.ts
|   |-- src/components/
|   |   |-- UploadModal.tsx
|   |   |-- QueryBox.tsx
|   |   |-- DocumentViewer.tsx
|   |   |-- EvidenceChain.tsx
|   |   `-- ClauseCard.tsx
|   `-- src/data/mockContract.ts
|-- backend-api/
|   |-- server.js
|   |-- db.js
|   |-- schema.sql
|   |-- services/parserService.js
|   |-- services/graphService.js
|   |-- services/geminiService.js
|   |-- services/queueService.js
|   |-- test_api.js
|   `-- test_e2e_integration.js
|-- backend-parser/
|   |-- main.py
|   |-- parser.py
|   |-- models.py
|   |-- test_service.py
|   `-- generate_demo_pdf.py
|-- report-latex/
|-- README.md
|-- RUN_DEMO.md
`-- VIVA_GUIDE.md
\end{lstlisting}

\subsection{What each major file does}
\begin{longtable}{>{\raggedright\arraybackslash}p{0.34\textwidth}>{\raggedright\arraybackslash}p{0.55\textwidth}}
\toprule
File & Responsibility\\
\midrule
\endfirsthead
\toprule
File & Responsibility\\
\midrule
\endhead
\path{frontend/src/App.tsx} & Controls main application state, selected document, mode and view composition.\\
\path{frontend/src/services/api.ts} & Health checks, live upload, live query, timeout handling and explicit offline behavior.\\
\path{frontend/src/components/EvidenceChain.tsx} & Displays retrieval, traversal, verification and verdict stages.\\
\path{frontend/src/data/mockContract.ts} & Static presentation data used only when sample/offline mode is deliberately selected.\\
\path{backend-api/server.js} & Express routes, PDF validation, orchestration and generation/fallback decisions.\\
\path{backend-api/services/parserService.js} & Sends uploaded bytes to the Python parser service.\\
\path{backend-api/services/graphService.js} & Stores clauses, resolves dependencies, scores queries, traverses the graph and detects missing evidence.\\
\path{backend-api/services/geminiService.js} & Builds the evidence-only prompt, calls Gemini, parses JSON and validates citation IDs.\\
\path{backend-api/db.js} & PostgreSQL access plus the development in-memory fallback.\\
\path{backend-api/schema.sql} & Database tables for contracts, clauses and clause dependencies.\\
\path{backend-parser/main.py} & FastAPI routes and request handling.\\
\path{backend-parser/parser.py} & Heading detection, clause construction, page tracking and reference extraction.\\
\path{backend-parser/models.py} & Typed request/response data models.\\
\path{backend-parser/test_service.py} & Parser tests and synthetic legal PDF generation.\\
\path{backend-api/test_e2e_integration.js} & Starts services, uploads a generated PDF, queries it and verifies missing-evidence abstention.\\
\bottomrule
\end{longtable}

\newpage
\subsection{Important APIs}
\begin{tabularx}{\linewidth}{>{\bfseries\raggedright\arraybackslash}p{0.28\linewidth}>{\raggedright\arraybackslash}X}
\toprule
Endpoint & Purpose\\
\midrule
GET /health & Reports Node status, database or in-memory mode, Python reachability, queue availability and Gemini/fallback status.\\
POST /api/contracts/upload & Accepts one PDF, invokes the parser, stores clauses, resolves dependencies and returns ingestion metadata.\\
POST /api/query & Runs seed retrieval, graph traversal, missing-evidence verification and optional answer generation.\\
GET /api/contracts & Lists stored contracts.\\
GET /api/contracts/:id & Returns one contract with clauses and dependency rows.\\
POST /parse-pdf & FastAPI endpoint returning metadata and structured clauses from PDF bytes.\\
\bottomrule
\end{tabularx}

\subsection{Core data objects}
\begin{lstlisting}
{
  "clause_id": "SEC-1.2",
  "section_title": "Exclusions from Protection",
  "page_number": 1,
  "content": "The obligations of Section 1.1 ...",
  "cross_references": ["Section 1.1", "Schedule A"]
}
\end{lstlisting}
A dependency records the contract, source clause, optional resolved target, original target code and relation type. Relation types include cross-reference, attachment reference, definition dependency and conditioned-on.

\newpage
\subsection{Testing vocabulary from zero}
Testing means giving a system controlled inputs and checking whether the actual result matches the expected result. It is not simply opening the application and observing that it looks correct.

\begin{longtable}{>{\bfseries\raggedright\arraybackslash}p{0.20\textwidth}>{\raggedright\arraybackslash}p{0.34\textwidth}>{\raggedright\arraybackslash}p{0.35\textwidth}}
\toprule
Term & Plain meaning & Example in this project\\
\midrule
\endfirsthead
\toprule
Term & Plain meaning & Example in this project\\
\midrule
\endhead
Test case & One defined situation containing an input, an action and an expected result. & Upload the synthetic two-page contract and expect a successful response containing extracted clauses.\\
Input & Data supplied to the code. & PDF bytes, a contract ID, or the question ``Is Schedule B attached?''\\
Expected result & The result decided before the test is run. & Schedule B is absent, so the query must abstain and report \texttt{SCHED-B}.\\
Actual result & What the program really returns. & The HTTP status and JSON body received from the API.\\
Assertion & A program statement that automatically compares the actual value with the expected value. & The code uses an assertion equivalent to ``missing result must have abstain equal to true.''\\
Pass & Every required assertion is true. & The parser found the required clause IDs and references.\\
Fail & At least one assertion is false or the program crashes. & The query returned an answer even though Schedule B was absent.\\
Bug or defect & A difference between required and actual behaviour. & A ``Section 1.1'' reference is normalized to the wrong clause ID.\\
Test fixture & Controlled data prepared for repeatable testing. & The generated NDA containing Sections 1.1, 1.2, 2.1, Clause 3.1 and Schedule A.\\
Synthetic data & Artificial but realistic data created for tests. & The contract is generated by code, so its correct structure is already known.\\
Deterministic & The same input and code should produce the same relevant output. & The missing-reference verifier returns the same refusal without depending on Gemini wording.\\
Positive test & Checks valid input and successful behaviour. & Upload a valid PDF and expect HTTP 201 from Node.\\
Negative test & Checks invalid input or a failure condition. & Upload a text file or an empty PDF and expect HTTP 400.\\
Regression test & A test preserved so that a future edit cannot silently break a working feature. & Keep the Schedule B refusal test after semantic retrieval is added.\\
\bottomrule
\end{longtable}

\subsection{The main levels of testing}
\begin{enumerate}
  \item \textbf{Unit testing} checks one small function or module in isolation. It is fast and helps locate a fault. In this project, reference extraction and reference normalization are tested directly.
  \item \textbf{API testing} sends requests to an endpoint and verifies status codes and JSON fields. Here, the FastAPI parser and Express routes are tested through HTTP-like requests.
  \item \textbf{Integration testing} checks whether multiple modules work together. Saving clauses and then resolving their dependency edges is an integration of storage and graph logic.
  \item \textbf{End-to-end or E2E testing} follows the complete user path across services. Our test generates a PDF, starts FastAPI and Express, uploads the PDF, performs queries and verifies the returned evidence.
  \item \textbf{Frontend static checking} uses TypeScript compilation and linting to detect type errors, unused values and suspicious code before the browser runs it. This is useful but is not the same as a complete browser test.
  \item \textbf{Evaluation} measures quality on many labelled examples using metrics. Evaluation is different from software testing: a test checks required behaviour, while evaluation estimates how accurate the research method is across a dataset.
\end{enumerate}

\subsection{Why several levels are necessary}
If a complete upload fails, the cause could be the PDF parser, the Node-to-Python request, storage, graph resolution or the endpoint. Unit tests narrow the location of the problem. Integration tests catch incorrect connections between modules. E2E testing gives confidence that the path visible to the user actually works. None of these alone proves research accuracy on unseen contracts; that requires the later evaluation dataset.

\newpage
\subsection{How testing is implemented in the current repository}
\begin{longtable}{>{\raggedright\arraybackslash}p{0.27\textwidth}>{\raggedright\arraybackslash}p{0.28\textwidth}>{\raggedright\arraybackslash}p{0.35\textwidth}}
\toprule
Location & What is exercised & What is checked\\
\midrule
\endfirsthead
\toprule
Location & What is exercised & What is checked\\
\midrule
\endhead
\path{backend-parser/test_service.py} & Python reference extraction & Schedule, Exhibit, Section, Clause and Article references are found by the extraction rules.\\
\path{backend-parser/test_service.py} & Structural PDF parsing & Required clause IDs, page metadata and cross-references are present after parsing the synthetic PDF.\\
\path{backend-parser/test_service.py} & FastAPI parser endpoint & Health succeeds; a valid PDF is parsed; a non-PDF and an empty PDF are rejected with status 400.\\
\path{backend-api/test_api.js} & Reference normalization & Human forms such as ``Section 1.1'' and ``Schedule A'' become stable IDs such as \texttt{SEC-1.1} and \texttt{SCHED-A}.\\
\path{backend-api/test_api.js} & Storage and dependency creation & Five mock clauses are stored, linked references resolve, and the contract receives an ID.\\
\path{backend-api/test_api.js} & Multi-hop graph retrieval & The exclusions seed is \texttt{SEC-1.2}; traversal reaches \texttt{SEC-1.1} and \texttt{SCHED-A}.\\
\path{backend-api/test_api.js} & Missing-evidence refusal & A question about absent Exhibit B produces \texttt{is\_abstain=true} and reports \texttt{EXHIBIT-B}.\\
\path{backend-api/test_api.js} & Query and listing routes & Query JSON, contract listing, contract details and evidence citation membership are checked.\\
\path{backend-api/test_e2e_integration.js} & Complete cross-service flow & A real generated PDF travels through Python and Node; upload, retrieval and missing Schedule B refusal are asserted.\\
\path{frontend/package.json} & Build and lint checks & \texttt{tsc -b}, Vite build and Oxlint check frontend types and source quality.\\
\bottomrule
\end{longtable}

\subsection{Commands used to conduct the current checks}
Run each command from the indicated folder. The E2E check needs FastAPI port 8000 and test port 3050 to be available.
\begin{lstlisting}
# From backend-parser/
.venv/bin/python test_service.py

# From backend-api/
npm test
PYTHON_BIN=../backend-parser/.venv/bin/python npm run test:e2e

# From frontend/
npm run build
npm run lint
\end{lstlisting}
The virtual-environment Python is specified because that environment contains PyMuPDF and the parser dependencies. A port-in-use message is an environment conflict: it means another server already owns that port, not automatically that an assertion in the tested logic failed.

\subsection{How the synthetic legal document is constructed}
The test does not depend on an unknown external PDF. Python and PyMuPDF create a two-page contract in memory. It deliberately contains:
\begin{itemize}
  \item Section 1.1 defining confidential information;
  \item Section 1.2 referring to Section 1.1 and the existing Schedule A;
  \item Section 2.1 referring to Section 1.1 and the deliberately absent Schedule B;
  \item Section 2.2 referring to Clause 3.1;
  \item Clause 3.1 referring to Section 2.1 and Article 2;
  \item Schedule A as a real target node.
\end{itemize}
Because the authors know which references exist and which one is intentionally missing, the document has known expected answers. This makes failures reproducible and separates a program defect from uncertainty in a real contract.

\subsection{Understanding HTTP and JSON in these tests}
HTTP is the communication protocol used between the browser, Node and Python. A method describes the action: GET reads data and POST submits data. An endpoint is a named route such as \texttt{/api/query}. A status code summarizes the outcome: 200 means a successful request, 201 means a resource was created, and 400 means the request was invalid. JSON is the structured text format used for responses. Tests inspect both the status code and important JSON values; checking only one would be incomplete.

\newpage
\subsection{The complete end-to-end test, step by step}
\begin{enumerate}
  \item \textbf{Create controlled input.} Python generates the synthetic PDF and writes it as the test contract.
  \item \textbf{Prepare the parser.} The test checks FastAPI on port 8000. If it is not already running, it starts Uvicorn with the parser application.
  \item \textbf{Verify parser health.} A health request must return status 200. This distinguishes an unavailable service from a parsing failure.
  \item \textbf{Start the API orchestrator.} Express starts on a separate test port so it does not conflict with the normal development server.
  \item \textbf{Verify cross-service connectivity.} Express health must report that the Python parser is reachable.
  \item \textbf{Upload through the public Node route.} The PDF is placed in multipart form data and sent to the contract-upload endpoint, \path{POST /api/contracts/upload}, exactly as a client uploads a file.
  \item \textbf{Assert ingestion results.} The response must be HTTP 201, successful, contain a contract ID, report at least six clauses and contain at least two linked dependencies.
  \item \textbf{Ask a resolvable question.} The test sends ``exclusions from protection'' with one seed and two hops. It requires a direct match and a grounded-answer envelope.
  \item \textbf{Ask a missing-evidence question.} The test asks whether Schedule B is attached and what it requires. Because the reference exists but the target does not, the correct behaviour is refusal.
  \item \textbf{Assert safe refusal.} The JSON must contain \texttt{is\_abstain=true}, provider \texttt{deterministic-verifier}, and normalized missing reference \texttt{SCHED-B}.
  \item \textbf{Clean up.} The test closes Express and terminates the Python process that it started.
\end{enumerate}

\subsection{Why assertions are stronger than visual observation}
The interface may display an attractive answer even when the wrong clause was retrieved. Automated assertions inspect internal facts: exact clause IDs, hop results, missing-reference codes, response provider and citation membership. If a later code change violates one of these expectations, the test exits with failure instead of relying on a person to notice it.

\subsection{Resolved and unresolved reference test}
Schedule A is present, so its reference should resolve to a graph node and should not trigger refusal. Schedule B is mentioned but absent, so its dependency keeps the normalized code with no target node. Testing both cases is important: always refusing would be as incorrect as never refusing.

\subsection{Gemini and deterministic testing}
Generative wording can vary, so the strongest current tests do not require one exact sentence from Gemini. They test deterministic controls around it: the evidence exists, the response has an answer envelope, every citation belongs to retrieved evidence, and missing evidence prevents generation. If no Gemini key is configured, an extractive or safe fallback still allows retrieval and verification to be tested.

\newpage
\subsection{Testing terms connected to research metrics}
\begin{longtable}{>{\bfseries\raggedright\arraybackslash}p{0.21\textwidth}>{\raggedright\arraybackslash}p{0.69\textwidth}}
\toprule
Term & Meaning for this project\\
\midrule
True positive & A required clause or missing-reference warning is correctly returned.\\
False positive & The system returns an irrelevant clause or wrongly claims that evidence is missing.\\
False negative & A required clause or genuinely missing reference is not detected.\\
Precision & Of everything returned, the proportion that was correct. High precision means little irrelevant evidence.\\
Recall & Of everything that should have been returned, the proportion the system found. High recall means fewer required clauses were missed.\\
F1 score & A balance of precision and recall, useful when both irrelevant and missing evidence matter.\\
Retrieval recall@$k$ & Whether the necessary evidence appears within the top $k$ initially retrieved clauses.\\
Citation validity & Whether every cited clause ID exists in the evidence supplied to the answer generator.\\
Citation completeness & Whether all important claims have supporting citations, not merely whether the IDs are real.\\
Abstention precision & When the system refuses, how often evidence really is insufficient.\\
Abstention recall & Of all cases that should be refused, how many the verifier catches.\\
False-refusal rate & How often the system refuses even though enough evidence exists.\\
Latency & Time taken for ingestion or a query. Accuracy is not useful if interactive review becomes unreasonably slow.\\
Baseline & A simpler method used for fair comparison. The proposed comparison is fixed-size retrieval versus structural dependency-aware retrieval.\\
Gold evidence & Clauses labelled by a human as necessary for answering a question; these labels are the reference for calculating metrics.\\
\bottomrule
\end{longtable}

\subsection{What has been conducted at the current stage}
At the present prototype stage, the team has conducted controlled functional tests rather than a final large-scale experiment. The implemented checks cover structural parsing, cross-reference extraction, valid and invalid file handling, reference normalization, clause storage, dependency creation, lexical seed retrieval, graph expansion, API responses, citation membership and missing-evidence abstention. A generated two-page contract is used so expected outputs are known.

The frontend is checked by TypeScript compilation, production build and linting. Manual demonstration checks confirm mode selection, upload feedback, evidence-chain rendering and visible refusal behaviour. Fully automated browser interaction tests have not yet been implemented.

\subsection{Representative end-to-end assertions}
The E2E test requires HTTP 201 for upload, a returned contract ID, at least six clauses and two dependencies in the synthetic contract, a direct match for the exclusions query, and a grounded-answer envelope. For the Schedule B question it requires \texttt{is\_abstain=true}, provider \texttt{deterministic-verifier}, and missing reference \texttt{SCHED-B}.

\subsection{Meaning of the current results}
The prototype validates feasibility, not final accuracy. It proves that real PDF bytes can travel through the services, become structured clauses, form dependency links, support a multi-hop query and trigger a genuine missing-evidence refusal. It does not yet prove performance on a large or diverse legal corpus.

\subsection{Testing still required}
\begin{itemize}
  \item Create a labelled dataset containing different contract types and layouts.
  \item Compare fixed-token, lexical, dense and hybrid retrieval with the same questions.
  \item Add scanned documents, OCR errors, tables, multi-column pages and malformed PDFs.
  \item Add automated browser tests for upload, query, mode switching and error messages.
  \item Test prompt injection hidden inside contract text and verify that it cannot override system rules.
  \item Measure precision, recall, citation completeness, refusal errors and latency.
  \item Repeat runs after every major change as regression tests.
\end{itemize}

\subsection{Additional material Aniket can explain}
Aniket can connect testing to reliability: a legal system must show not only an answer but how the evidence travelled through the system. He can explain the difference between a unit test and an E2E test, why synthetic data gives known answers, why both Schedule A and missing Schedule B are necessary, why deterministic verification is easier to test than free-form LLM text, and why present passing tests demonstrate engineering feasibility rather than universal accuracy.

\subsection{Current limitations}
\begin{itemize}
  \item Rule-based structure detection depends on recognizable headings.
  \item Scanned PDFs require OCR, which is not yet integrated.
  \item Block reading order may be imperfect in complex layouts.
  \item Lexical retrieval suffers from vocabulary mismatch.
  \item The main prototype reasons over one contract at a time.
  \item In-memory storage disappears when the process restarts.
  \item A language model can still paraphrase evidence incorrectly.
  \item Privacy, authentication and deployment hardening are not final.
\end{itemize}

\subsection{Security and privacy theory}
Keep the Gemini key only in backend environment variables; validate file type and size; treat contract text as untrusted data; send only necessary evidence; add authenticated workspaces and authorization; define retention and deletion; encrypt data in transit and at rest; record audit events; and test prompt injection inside contract text.

\subsection{Remaining roadmap}
Improve parsing and OCR; build an evaluation collection; add dense embeddings, BM25 and reranking; measure retrieval and refusal behavior; model amendments and supersession; use persistent PostgreSQL/pgvector; add authentication and privacy controls; conduct adversarial and usability tests; then deploy and write the final comparative report.

\newpage
\section{Shared Technical Questions and Short Answers}
\textbf{Is this RAG?} It is RAG-related. The current baseline uses clause-aware lexical retrieval, graph expansion and verification; dense vector retrieval is planned.\par
\textbf{Why not send the entire PDF to Gemini?} Selecting evidence first controls context, cost, distraction and auditability.\par
\textbf{Does the system train Gemini?} No. Gemini is pretrained; the project controls its evidence and validates citations.\par
\textbf{Why is the graph directed?} The direction records which source clause explicitly depends on which target.\par
\textbf{Why two hops?} It captures nearby chains while controlling context growth and latency; the best depth should later be evaluated.\par
\textbf{How is a missing schedule detected?} An unresolved dependency has a reference code but no target node. If relevant to the query, it triggers abstention.\par
\textbf{What is hard-coded?} Sample/offline data and baseline rules are coded. Live document content is genuinely uploaded and parsed.\par
\textbf{Can it process every legal PDF?} No. Text-based consistent contracts work best; scans and complex layouts need future improvements.\par
\textbf{What if Gemini is unavailable?} Verified evidence is returned through the fallback. Retrieval and abstention do not depend on Gemini.\par
\textbf{What prevents a fake citation?} Returned citation IDs are checked against the evidence allow-list.\par
\textbf{Why separate Python and Node?} Python provides PDF tooling; Node fits web orchestration and React integration.\par
\textbf{What is the final research comparison?} Structural dependency-verified retrieval versus fixed-size chunking on evidence, citations, correctness, refusal and latency.\par

\section{Study Routine Without Speech Memorization}
\subsection{Build the common mental model}
All members draw the pipeline from memory: PDF $\rightarrow$ parser $\rightarrow$ clauses $\rightarrow$ dependency graph $\rightarrow$ seed retrieval $\rightarrow$ graph expansion $\rightarrow$ verifier $\rightarrow$ Gemini/fallback $\rightarrow$ evidence UI. Explain every arrow.

\subsection{Learn assigned material}
\begin{itemize}
  \item Puru studies the introduction, layers, both data flows, scoring, graph, Gemini control and technology choices.
  \item Sujal studies the four papers, their differences from this project, five gaps, hypotheses and metrics.
  \item Aniket studies the repository, files, endpoints, data objects, tests, results, limitations, security and roadmap.
\end{itemize}

\subsection{Explain without reading}
Puru redraws the architecture. Sujal redraws the gap-to-solution mapping. Aniket redraws the folder tree and test pipeline. Compare with this guide and repeat missing parts.

\subsection{Cross-questioning and demo verification}
Each member asks another member questions from both shared and assigned sections. Then all members watch one live upload and both demonstration questions. Everyone must identify the seed clause, graph-expanded clauses, unresolved Schedule B edge, abstention stage and answer provider.

\subsection{Final accuracy check}
Never claim 100 percent accuracy, complete hallucination elimination, full security, universal PDF support, a trained proprietary model or completed multi-document reasoning. The accurate claim is that the prototype validates an end-to-end evidence-first foundation and defines a measurable research plan.

\section{References}
\begin{enumerate}
  \item P. Lewis et al., ``Retrieval-Augmented Generation for Knowledge-Intensive NLP Tasks,'' \textit{NeurIPS}, 2020. \url{https://arxiv.org/abs/2005.11401}
  \item D. Hendrycks, C. Burns, A. Chen, and S. Ball, ``CUAD: An Expert-Annotated NLP Dataset for Legal Contract Review,'' 2021. \url{https://arxiv.org/abs/2103.06268}
  \item N. Guha et al., ``LegalBench: A Collaboratively Built Benchmark for Measuring Legal Reasoning in Large Language Models,'' 2023. \url{https://arxiv.org/abs/2308.11462}
  \item D. Edge et al., ``From Local to Global: A Graph RAG Approach to Query-Focused Summarization,'' 2024. \url{https://arxiv.org/abs/2404.16130}
\end{enumerate}
\end{document}
